\subsection{LU, QR, and Cholesky-Type Decompositions}

Throughout this subsection, $\mathbb{F}$ is a field.

\begin{definition} \label{def:full-rank-factorization} \hlt{Full-Rank Factorization}\\
Let $m,n,r\in\mathbb{Z}_{>0}$ and $A\in\operatorname{Mat}(m,n,\mathbb{F})$.\\
A factorization $A=XY$, with $X\in\operatorname{Mat}(m,r,\mathbb{F})$ and $Y\in\operatorname{Mat}(r,n,\mathbb{F})$, is a \hlt{full-rank factorization} if $\operatorname{rank}(X)=\operatorname{rank}(Y)=r$.
\end{definition}

\begin{theorem} \label{thm:full-rank-factorization} \hlt{Full-Rank Factorization}\\
Let $m,n\in\mathbb{Z}_{>0}$ and let $A\in\operatorname{Mat}(m,n,\mathbb{F})$ have rank $r\geq1$. Let the columns of $X\in\operatorname{Mat}(m,r,\mathbb{F})$ form a basis of $\operatorname{col}(A)$. Then:
\begin{enumerate}[label=(\roman*)]
\setlength{\itemsep}{0pt}
\item $(\exists!\,Y\in\operatorname{Mat}(r,n,\mathbb{F}))\,A=XY$;
\item for this matrix $Y$, $\operatorname{rank}(X)=\operatorname{rank}(Y)=r$.
\end{enumerate}
\begin{proof}
\begin{enumerate}[label=(\roman*)]
\setlength{\itemsep}{0pt}
\item Write $A=[\mathbf{a}_1\ \cdots\ \mathbf{a}_n]$. By the basis hypothesis, $(\forall\,j\in J_n)(\exists!\,\mathbf{y}_j\in\mathbb{F}^r)\,\mathbf{a}_j=X\mathbf{y}_j$.\\
For $Y\defeq[\mathbf{y}_1\ \cdots\ \mathbf{y}_n]$, one has $A=XY$. If $A=XZ$ with $Z=[\mathbf{z}_1\ \cdots\ \mathbf{z}_n]$, then $Z=Y$ by coordinate uniqueness.
\item $\operatorname{rank}(X)=r$ by the basis hypothesis, and $r=\operatorname{rank}(XY)\leq\operatorname{rank}(Y)\leq r$ by Theorem \ref{thm:rank-product-and-concatenation-bounds}(i) and the row count of $Y$.
\end{enumerate}
\hlt{Structural Commentary}: Basis coordinates construct $Y$ uniquely; the product-rank bound forces full row rank.\\
\hlt{Key Insight}: Column coordinates determine the right factor.
\end{proof}
\end{theorem}

\begin{definition} \label{def:lu-factorization} \hlt{LU Factorisation}\\
Let $n\in\mathbb{Z}_{>0}$ and $A\in\operatorname{Mat}(n,n,\mathbb{F})$.
\begin{enumerate}[label=(\roman*)]
\setlength{\itemsep}{0pt}
\item An \hlt{LU factorisation} is a product $A=LU$ with $L$ lower triangular and $U$ upper triangular.
\item For $L=[l_{ij}]$, it is a \hlt{unit-lower-triangular LU factorisation} if $(\forall\,k\in J_n)\,l_{kk}=1$.
\end{enumerate}
\end{definition}

\begin{theorem} \label{thm:lu-factorization-leading-principal-minors} \hlt{LU Factorisation from Nonzero Leading Principal Minors}\\
Let $n\in\mathbb{Z}_{>0}$ and $A\in\operatorname{Mat}(n,n,\mathbb{F})$. Assume every leading principal minor of $A$ is nonzero.\\
Then a unique pair $(L,U)\in\operatorname{Mat}(n,n,\mathbb{F})^2$ satisfies $A=LU$, with $L$ unit lower triangular and $U$ upper triangular.
\begin{proof}
Induct on $n$.
\begin{enumerate}[label=(\roman*)]
\setlength{\itemsep}{0pt}
\item \hlt{Base Case}: If $n=1$, then $A=[a]=[1][a]$ uniquely, where $a\neq0$ by hypothesis.
\item \hlt{Schur Reduction}: For $n\geq2$, write $A=\left[\begin{smallmatrix}a&\mathbf{r}\\\mathbf{c}&B\end{smallmatrix}\right]$ with $a\neq0$, and set $\boldsymbol{\ell}\defeq a^{-1}\mathbf{c}$ and $S\defeq B-a^{-1}\mathbf{c}\mathbf{r}$.\\
For the leading blocks $A_k$ and $S_{k-1}$, $\det(A_k)=a\det(S_{k-1})$ by Proposition \ref{prop:schur-complement-determinant-formula}.\\
Every leading principal minor of $S$ is nonzero by this identity and the hypothesis. Hence $S=L_1U_1$ uniquely, with $L_1$ unit lower triangular and $U_1$ upper triangular, by the induction hypothesis.
\item \hlt{Construction}: Set $L\defeq\left[\begin{smallmatrix}1&0\\\boldsymbol{\ell}&L_1\end{smallmatrix}\right]$ and $U\defeq\left[\begin{smallmatrix}a&\mathbf{r}\\0&U_1\end{smallmatrix}\right]$. Then $L$ and $U$ have the required forms, and $LU=\left[\begin{smallmatrix}a&\mathbf{r}\\\mathbf{c}&B\end{smallmatrix}\right]=A$.
\item \hlt{Uniqueness}: For another admissible pair, write $\widetilde L=\left[\begin{smallmatrix}1&0\\\widetilde{\boldsymbol\ell}&\widetilde L_1\end{smallmatrix}\right]$ and $\widetilde U=\left[\begin{smallmatrix}\widetilde a&\widetilde{\mathbf r}\\0&\widetilde U_1\end{smallmatrix}\right]$.\\
Then $(\widetilde a,\widetilde{\mathbf r},\widetilde{\boldsymbol\ell})=(a,\mathbf r,\boldsymbol\ell)$ by the first row and column. Then $S=\widetilde L_1\widetilde U_1$, so $(\widetilde L_1,\widetilde U_1)=(L_1,U_1)$ by the induction hypothesis.
\end{enumerate}
\hlt{Structural Commentary}: Schur reduction transfers the pivot hypothesis before block reconstruction and uniqueness.\\
\hlt{Key Insight}: The first pivot fixes the border blocks.
\end{proof}
\end{theorem}

\begin{corollary} \label{cor:symmetric-lu-gives-ldlt} \hlt{Symmetric LU and an $LDL^t$ Factorisation}\\
Let $n\in\mathbb{Z}_{>0}$ and let $A=A^t\in\operatorname{Mat}(n,n,\mathbb{R})$ have nonzero leading principal minors.\\
Then $(\exists\,L,D\in\operatorname{Mat}(n,n,\mathbb{R}))\,A=LDL^t$, where $L$ is unit lower triangular and $D$ is diagonal and invertible.
\begin{proof}
By Theorem \ref{thm:lu-factorization-leading-principal-minors}, write $A=LU$ uniquely with $L$ unit lower triangular and $U$ upper triangular.\\
The relation $0\neq\det(A)=\det(U)=\prod\limits_{k=1}^{n}u_{kk}$ holds by Theorem \ref{thm:determinant-basic-laws}(v)--(vi) and $\det(L)=1$.\\
Thus $D\defeq\operatorname{diag}(u_{11},\ldots,u_{nn})$ is invertible. For $U_0\defeq D^{-1}U$, one has $A=LDU_0$, with $U_0$ unit upper triangular.\\
By symmetry, $A=U_0^t(DL^t)$ is another admissible LU factorisation; hence $U_0^t=L$ by Theorem \ref{thm:lu-factorization-leading-principal-minors}. Thus $A=LDL^t$.\\
\hlt{Structural Commentary}: Pivot extraction normalises $U$; symmetry and LU uniqueness identify its unit factor.\\
\hlt{Key Insight}: Symmetry forces $U_0=L^t$.
\end{proof}
\end{corollary}

\begin{definition} \label{def:wide-and-narrow-qr-factorization} \hlt{Wide and Narrow QR Factorisation}\\
Let $m,n\in\mathbb{Z}_{>0}$ satisfy $m\geq n$ and let $\mathbb{F}\in\{\mathbb{R},\mathbb{C}\}$. Let $A\in\operatorname{Mat}(m,n,\mathbb{F})$.
\begin{enumerate}[label=(\roman*)]
\setlength{\itemsep}{0pt}
\item A \hlt{wide QR factorisation} is $A=V\left[\begin{smallmatrix}R\\0\end{smallmatrix}\right]$, where $V\in\operatorname{Mat}(m,m,\mathbb{F})$ is unitary and $R\in\operatorname{Mat}(n,n,\mathbb{F})$ is upper triangular.
\item For $V=[Q\ Q']$ with $Q\in\operatorname{Mat}(m,n,\mathbb{F})$, the corresponding \hlt{narrow QR factorisation} is $A=QR$, where $Q^*Q=I_n$.
\end{enumerate}
\end{definition}

\begin{theorem} \label{thm:qr-factorization} \hlt{QR Factorisation}\\
Let $m,n\in\mathbb{Z}_{>0}$ satisfy $m\geq n$ and let $\mathbb{F}\in\{\mathbb{R},\mathbb{C}\}$. Let $A\in\operatorname{Mat}(m,n,\mathbb{F})$.
\begin{enumerate}[label=(\roman*)]
\setlength{\itemsep}{0pt}
\item $(\exists\,V\in\operatorname{Mat}(m,m,\mathbb{F}))(\exists\,R\in\operatorname{Mat}(n,n,\mathbb{F}))\,A=V\left[\begin{smallmatrix}R\\0\end{smallmatrix}\right]$. Here $V$ is unitary, $R=[r_{ij}]$ is upper triangular with $(\forall\,k\in J_n)\,r_{kk}\in\mathbb{R}_{\geq0}$.\\
For $V=[Q\ Q']$ with $Q\in\operatorname{Mat}(m,n,\mathbb{F})$, one has $A=QR$ and $Q^*Q=I_n$.
\item If $\operatorname{rank}(A)=n$, then a unique pair $(Q,R)$ satisfies $A=QR$ and $Q^*Q=I_n$.\\
Here $Q\in\operatorname{Mat}(m,n,\mathbb{F})$ and $R=[r_{ij}]\in\operatorname{Mat}(n,n,\mathbb{F})$ is upper triangular with $(\forall\,k\in J_n)\,r_{kk}>0$.
\end{enumerate}
\begin{proof}
\begin{enumerate}[label=(\roman*)]
\setlength{\itemsep}{0pt}
\item Process $k\in J_n$ in increasing order, and let $\mathbf{x}_k\in\mathbb{F}^{m-k+1}$ be the active tail of column $k$.\\
For $\mathbf{x}_k=\mathbf{0}$, set $(H_k,r_{kk})\defeq(I_{m-k+1},0)$. Otherwise set $\mathbf{y}_k\defeq\norm{\mathbf{x}_k}_2\mathbf{e}_1$ and choose $\sigma_k,\mathbf{w}_k$ by Theorem \ref{thm:householder-maps-equal-norm-vectors}(i)--(iii).\\
Then $H_k\defeq\sigma_kU_{\mathbf{w}_k}$ satisfies $H_k^*H_k=I_{m-k+1}$ by Theorem \ref{thm:householder-maps-equal-norm-vectors}(i) and Theorem \ref{thm:householder-matrix-properties}(iii).\\
It also satisfies $H_k\mathbf{x}_k=r_{kk}\mathbf{e}_1$ for $r_{kk}\defeq\norm{\mathbf{x}_k}_2$ by Theorem \ref{thm:householder-maps-equal-norm-vectors}(iii).\\
In both cases, $\widehat H_k\defeq I_{k-1}\oplus H_k$ preserves prior zeros by acting only on rows $k,\ldots,m$.\\
Thus $WA=\left[\begin{smallmatrix}R\\0\end{smallmatrix}\right]$ for $W\defeq\widehat H_n\cdots\widehat H_1$, with $R$ upper triangular and $(\forall\,k\in J_n)\,r_{kk}\geq0$.\\
For $V\defeq W^*=[Q\ Q']$, $A=V\left[\begin{smallmatrix}R\\0\end{smallmatrix}\right]=QR$ and $Q^*Q=I_n$ by unitarity and block multiplication.
\item Assume $\operatorname{rank}(A)=n$. Then $\operatorname{rank}(R)=n$ by Theorem \ref{thm:full-rank-products-preserve-rank}(ii), so $R$ is invertible by Corollary \ref{cor:square-matrix-invertible-iff-full-rank}.\\
Its nonnegative diagonal is positive by Theorem \ref{thm:determinant-basic-laws}(v),(vii), proving existence.\\
If $A=Q_1R_1=Q_2R_2$ are admissible, then each $R_i$ is invertible by Theorem \ref{thm:determinant-basic-laws}(v),(vii), and $R_1^*R_1=A^*A=R_2^*R_2$.\\
For $D\defeq R_1R_2^{-1}$, $D$ is upper triangular by Theorem \ref{thm:inverse-upper-triangular}(ii) and the product-entry formula.\\
The identity $D^*D=(R_2^{-1})^*R_1^*R_1R_2^{-1}=I_n$ holds, and $DD^*=I_n$ by Theorem \ref{thm:unitary-matrix-equivalent-conditions}(i)$\Rightarrow$(iii).\\
Viewed over $\mathbb{C}$ when $\mathbb{F}=\mathbb{R}$, $D$ is diagonal by Lemma \ref{lem:upper-triangular-normal-matrix-is-diagonal}.\\
Now $d_{kk}=(R_1)_{kk}(R_2)_{kk}^{-1}>0$ and $\abs{d_{kk}}=1$, so $D=I_n$. Therefore $R_1=R_2$ and $Q_1=AR_1^{-1}=Q_2$.
\end{enumerate}
\hlt{Structural Commentary}: Householder tail maps construct $R$; triangular normality reduces uniqueness to diagonal phases.\\
\hlt{Key Insight}: Positivity eliminates every unitary diagonal phase.
\end{proof}
\end{theorem}

\begin{corollary} \label{cor:qr-prefix-column-spaces} \hlt{Prefix Column Spaces in QR Factorisation}\\
Let $m,n\in\mathbb{Z}_{>0}$ satisfy $m\geq n$ and let $\mathbb{F}\in\{\mathbb{R},\mathbb{C}\}$.\\
Let $A=QR\in\operatorname{Mat}(m,n,\mathbb{F})$ satisfy $\operatorname{rank}(A)=n$ and $Q^*Q=I_n$, where $R$ is upper triangular with positive diagonal.\\
For $k\in J_n$, let $A_k$ and $Q_k$ be the first $k$ columns of $A$ and $Q$. Then:
\begin{enumerate}[label=(\roman*)]
\setlength{\itemsep}{0pt}
\item $(\forall\,k\in J_n)\,\operatorname{col}(A_k)=\operatorname{col}(Q_k)$;
\item $(\forall\,k\in J_n)$, the columns of $Q_k$ form an orthonormal basis of $\operatorname{col}(A_k)$.
\end{enumerate}
\begin{proof}
\begin{enumerate}[label=(\roman*)]
\setlength{\itemsep}{0pt}
\item Fix $k\in J_n$, and let $R_k$ be the leading $k\times k$ block of $R$. Then $R_k$ is invertible by Theorem \ref{thm:determinant-basic-laws}(v),(vii).\\
Since $A_k=Q_kR_k$ and $Q_k=A_kR_k^{-1}$, one has $\operatorname{col}(A_k)=\operatorname{col}(Q_k)$.
\item Fix $k\in J_n$. The columns of $Q_k$ are orthonormal by $Q_k^*Q_k=I_k$, and they span $\operatorname{col}(A_k)$ by part (i).
\end{enumerate}
\hlt{Structural Commentary}: Invertible triangular prefixes transfer column spaces; $Q_k^*Q_k=I_k$ supplies orthonormality.\\
\hlt{Key Insight}: Each prefix uses only the first $k$ orthonormal columns.
\end{proof}
\end{corollary}

Figure \ref{fig:qr-factorization-geometry} depicts orthonormal column replacement and upper-triangular coordinates.

\begin{illustrationfigure}[H]
  \centering
  \begin{tikzpicture}[
  >=Stealth,
  every node/.style={font=\scriptsize},
  box/.style={draw=gray!55, rounded corners=3pt, inner sep=5pt, align=center},
  arrow/.style={->, line width=0.85pt, blue!70!black},
  brace/.style={gray!60, line width=0.7pt}
]
  \node[box] (A) at (0,0)
    {$A=[\mathbf{a}_1\ \cdots\ \mathbf{a}_n]$\\original columns};
  \node[box] (Q) at (5.15,0)
    {$Q=[\mathbf{q}_1\ \cdots\ \mathbf{q}_n]$\\orthonormal columns};
  \node[box] (R) at (11.25,0)
    {$R$ upper triangular\\coordinates in the $Q$ basis};

  \draw[arrow] (A) -- (Q) node[midway, above=4pt, fill=white, inner sep=1pt] {orthonormalise};
  \draw[arrow] (Q) -- (R) node[midway, above=4pt, fill=white, inner sep=1pt] {record coefficients};
\end{tikzpicture}

  \caption{QR records the original columns in an orthonormal basis while preserving every prefix column space.}
  \label{fig:qr-factorization-geometry}
\end{illustrationfigure}

\begin{perspective} \label{per:gram-schmidt-householder-qr} \hlt{Gram--Schmidt and Householder QR}\\
For a full-column-rank matrix, Gram--Schmidt preserves prefix spans by Theorem \ref{thm:gram-schmidt-process}(i)--(ii).\\
Householder transformations instead triangularise entire column tails in Theorem \ref{thm:qr-factorization}(i).\\
Both constructions reach the unique positive-diagonal factorisation in Theorem \ref{thm:qr-factorization}(ii).\\
Gram--Schmidt orthogonalises columns; Householder transformations act unitarily on rows.
\end{perspective}

\begin{perspective} \label{per:rank-deficient-gram-schmidt-qr} \hlt{Rank-Deficient QR from Gram--Schmidt}\\
Zero Gram--Schmidt remainders detect dependent prefix columns by Corollary \ref{cor:gram-schmidt-zero-remainders}.\\
The QR factorisation still exists by Theorem \ref{thm:qr-factorization}(i).\\
Rank defect blocks the positive-diagonal uniqueness hypothesis in Theorem \ref{thm:qr-factorization}(ii).
\end{perspective}

\begin{theorem} \label{thm:cholesky-factorization} \hlt{Cholesky Factorisation}\\
Let $n\in\mathbb{Z}_{>0}$, let $\mathbb{F}\in\{\mathbb{R},\mathbb{C}\}$, and let $A\in\operatorname{Mat}(n,n,\mathbb{F})$.
\begin{enumerate}[label=(\roman*)]
\setlength{\itemsep}{0pt}
\item If $A\succeq0$, then $(\exists\,L\in\operatorname{Mat}(n,n,\mathbb{F}))\,A=LL^*$, where $L$ is lower triangular with nonnegative diagonal entries.
\item If $A\succ0$, then $\exists!\,L\in\operatorname{Mat}(n,n,\mathbb{F})$ satisfying $A=LL^*$, where $L$ is lower triangular with positive diagonal entries.
\end{enumerate}
\begin{proof}
\begin{enumerate}[label=(\roman*)]
\setlength{\itemsep}{0pt}
\item Set $B\defeq A^{1/2}$ by Corollary \ref{cor:positive-semidefinite-square-root-formula}(i), so $B=B^*$ and $B^2=A$.\\
Choose $B=QR$ with $Q^*Q=I_n$ and $R$ upper triangular with nonnegative diagonal by Theorem \ref{thm:qr-factorization}(i).\\
Then $A=B^*B=R^*R$, so $L\defeq R^*$ has the required form.
\item Assume $A\succ0$. Then $A$ is invertible by Theorem \ref{thm:positive-definite-matrix-equivalences}(i)$\Rightarrow$(iii).\\
For the factor $L$ in (i), $n=\operatorname{rank}(A)\leq\operatorname{rank}(L)\leq n$ by Theorem \ref{thm:rank-product-and-concatenation-bounds}(i).\\
Thus $L$ is invertible by Corollary \ref{cor:square-matrix-invertible-iff-full-rank}, so its nonnegative diagonal is positive by Theorem \ref{thm:determinant-basic-laws}(v),(vii).\\
For uniqueness, set $B\defeq A^{1/2}$ and suppose $A=L_1L_1^*=L_2L_2^*$ for admissible $L_1,L_2$.\\
Since $B^2=A$ and $A$ is invertible, $\ker(B)\subseteq\ker(A)=\{\mathbf{0}\}$; hence $B$ is invertible by Corollary \ref{cor:square-matrix-algebraic-invertibility-dictionary}(vi)$\Rightarrow$(i).\\
For $R_i\defeq L_i^*$, each $R_i$ is upper triangular with positive diagonal and is invertible by Theorem \ref{thm:determinant-basic-laws}(v),(vii). Set $Q_i\defeq BR_i^{-1}$.\\
Then $B=Q_iR_i$ and $Q_i^*Q_i=(R_i^{-1})^*B^2R_i^{-1}=I_n$. Hence $R_1=R_2$ by Theorem \ref{thm:qr-factorization}(ii), and therefore $L_1=L_2$.
\end{enumerate}
\hlt{Structural Commentary}: QR of $A^{1/2}$ constructs $R^*$; positive-diagonal QR uniqueness identifies competitors.\\
\hlt{Key Insight}: Cholesky is the adjoint triangular factor of $A^{1/2}$.
\end{proof}
\end{theorem}

Figure \ref{fig:cholesky-from-qr} depicts the Cholesky factor obtained from QR of the positive square root.

\begin{illustrationfigure}[H]
  \centering
  \begin{tikzpicture}[
  >=Stealth,
  x=0.86cm,
  every node/.style={font=\small},
  box/.style={draw=gray!55, rounded corners=3pt, fill=gray!6, align=center, inner sep=4pt},
  arrow/.style={->, line width=0.85pt, blue!70!black}
]
  \node[box] (root) at (0,0) {$B=A^{1/2}$};
  \node[box] (qr) at (2.7,0) {$B=QR$};
  \node[box] (chol) at (5.35,0) {$L=R^*$};

  \draw[arrow] (root) -- (qr) node[midway, above] {QR};
  \draw[arrow] (qr) -- (chol) node[midway, above] {adjoint};
\end{tikzpicture}

  \caption{For $B=A^{1/2}=QR$, the choice $L=R^*$ satisfies $A=LL^*$; positive definiteness makes $L$ unique.}
  \label{fig:cholesky-from-qr}
\end{illustrationfigure}
